\begin{textsample}{POS Dim 2 – es – Score 19.00 – es/es\_inf\_intro\_3.txt}  \label{ex:f2_pos_247}
A Elaboração do \textbf{Currículo} em Regime de Colaboração A definição de uma base comum curricular para todo o país atende a uma \textbf{prerrogativa} da Constituição Federal Brasileira ( BRASIL, 1988 ), da Lei de \textbf{Diretrizes} e Bases da Educação Nacional ( BRASIL, 1996 ) e do Plano Nacional de Educação 2014-2024 ( BRASIL, 2014 ) e nos coloca no rumo dos principais sistemas educacionais do mundo.
Ao mesmo tempo, nos desafia a ter um novo olhar sobre os \textbf{currículos} já construídos e vividos nas redes estaduais e \textbf{municipais} de ensino, pois passa a ser uma referência nacional obrigatória para elaboração ou revisão dos documentos curriculares.
A construção do \textbf{Currículo} do Espírito Santo se dá num momento histórico da educação brasileira.
Em 17 de dezembro de 2017 foi \textbf{homologada} pelo Conselho Nacional de Educação a Base Nacional Comum Curricular, para a Educação Infantil e o Ensino Fundamental, que estabelece as aprendizagens essenciais e indispensáveis a todos os estudantes da educação básica nessas etapas¹.
A definição de uma base comum curricular para todo o país atende a uma \textbf{prerrogativa} da Constituição Federal Brasileira de 1988, da Lei de \textbf{Diretrizes} e Bases da Educação Nacional ( Lei Nº 9394/96 ) e do Plano Nacional de Educação de 2014 e nos coloca no rumo dos principais sistemas educacionais do mundo.
Ao mesmo tempo, nos desafia a ter um novo olhar sobre os \textbf{currículos} já construídos e vividos nas redes estaduais e \textbf{municipais} de ensino, pois passa a ser uma referência nacional obrigatória para elaboração ou revisão curricular.
Nesse contexto, o Ministério da Educação \textbf{instituiu}, na \textbf{Portaria} Nº 331, de 5 de abril de 2018, o Programa de Apoio à Implementação da Base Nacional Comum Curricular – ProBNCC, cuja adesão pela Secretaria de Estado da Educação do Espírito Santo - SEDU e União Nacional dos \textbf{Dirigentes} \textbf{Municipais} de Educação, seccional Espírito Santo - UNDIME/ES, revela o compromisso das duas \textbf{instituições} em construir um \textbf{currículo}, em regime de colaboração entre estado e \textbf{municípios}, para proporcionar uma dinâmica de continuidade na formação do estudante de todo o território espírito-santense e desenvolver uma visão integrada para o desenvolvimento das ações necessárias para implementação e \textbf{gestão} curricular.
Para o desenvolvimento de um trabalho de tal magnitude, foi \textbf{instituída}, pela \textbf{Portaria} Nº037-R/2018, uma estrutura de governança, visando dar assento, em igualdade, a instâncias representativas do estado e \textbf{municípios}, bem como a \textbf{instituições} que representam os profissionais da educação e as que são responsáveis por sua formação.
Na mesma \textbf{portaria} foi \textbf{instituída} a \textbf{equipe} de elaboração curricular, composta por duas coordenações estaduais ( CONSED e UNDIME ), três coordenações estaduais de etapa ( Educação Infantil e Ensino Fundamental - Anos Iniciais e Anos Finais ), um analista de \textbf{gestão}, um articulador de regime de colaboração e 19 \textbf{redatores} dos componentes curriculares elencados na BNCC, além dos articuladores do Conselho Estadual de Educação - CEE e da União dos Conselhos \textbf{Municipais} de Educação - UNCME.
Importante mencionar que a \textbf{equipe} de \textbf{redatores} foi composta por professores das redes estadual e \textbf{municipal}, que convidaram outros professores colaboradores de diferentes redes para contribuir com a elaboração desse documento.
Além do estudo profundo da Base Nacional Comum Curricular, a \textbf{equipe} de \textbf{currículo} realizou estudos dos documentos normativos e legais da educação nacional ( Constituição Federal de 1988, LDB 9394/96, Diretrizes Nacionais da Educação Básica : Diversidade e Inclusão de 2013 ), de \textbf{currículos} nacionais e internacionais, e, principalmente, dos \textbf{currículos} já construídos e vividos na rede estadual, no caso o \textbf{Currículo} Básico Escola Estadual - CBEE ( ES, 2009 ), e nas redes \textbf{municipais} do Espírito Santo².
No seu processo de elaboração, o documento passou por duas consultas públicas online, a primeira direcionada aos profissionais de educação e a segunda também aberta para a sociedade ; bem como por leitura crítica de profissionais e \textbf{instituições} representativas que desenvolvem estudos e pesquisas, uma vez que influenciam na construção de \textbf{políticas} públicas e formação profissional de professores nas diversas áreas e etapas que são abrangidas pelo \textbf{currículo}.
Há que se destacar ainda o papel imprescindível dos articuladores \textbf{municipais}, indicados por suas \textbf{secretarias}, das SREs e professores referência, na mobilização dos professores e demais profissionais da educação de suas redes para que fossem protagonistas da construção coletiva e colaborativa deste documento curricular, que no total recebeu 10.649 contribuições de profissionais da educação e da sociedade civil.
O \textbf{Currículo} do Espírito Santo, construído por muitos sujeitos, é resultado do trabalho em conjunto entre as \textbf{instituições} parceiras e a \textbf{equipe} de \textbf{currículo} e da colaboração de diversos profissionais da educação dos mais diferentes lugares de nosso estado, o que permitiu o avanço das propostas inicialmente apresentadas e uma visão mais integrada do percurso formativo dos estudantes da educação básica de nosso território, que direcionará outras \textbf{políticas} e ações necessárias para a sua implementação nas \textbf{secretarias} e escolas estaduais e \textbf{municipais}, incluindo orientações didático-metodológicas, materiais didáticos e formação docente.
Importante destacar que o \textbf{Currículo} do Espírito Santo contempla os componentes curriculares abordados pela Base Nacional Comum Curricular, que define as aprendizagens essenciais dos componentes obrigatórios em todos os \textbf{currículos}, e os contextualiza, aprofunda e complementa nas questões relativas à educação do nosso Estado.
Cabe a cada rede, adesa a este documento, elaborar outros componentes que sejam exigidos por normas específicas ao seu contexto. ¹ Quando \textbf{homologadas} as aprendizagens essenciais do Ensino Médio, elas serão incorporadas a esse documento.

% matched lemmas: currículo, diretriz, dirigente, equipe, gestão, homologar, instituir, instituição, municipal, município, política, portaria, prerrogativa, redator, secretaria
\end{textsample}
